% Generated draft: requires mathematical and source review.
\documentclass[11pt]{article}
\usepackage[margin=25mm]{geometry}
\usepackage{fontspec}
\setmainfont{Latin Modern Roman}
\usepackage{amsmath,amssymb,mathtools,unicode-math}
\setmathfont{Latin Modern Math}
\usepackage{enumitem,needspace,xurl,hyperref}
\setlength{\parindent}{0pt}
\setlength{\parskip}{5pt}
\setlength{\emergencystretch}{3em}
\newcommand{\N}{\mathbb{N}}
\newcommand{\Z}{\mathbb{Z}}
\newcommand{\Q}{\mathbb{Q}}
\newcommand{\R}{\mathbb{R}}
\newcommand{\C}{\mathbb{C}}
\newcommand{\F}{\mathbb{F}}
\newcommand{\A}{\mathbb{A}}
\newcommand{\PP}{\mathbb{P}}
\newcommand{\E}{\mathbb{E}}
\newcommand{\eps}{\varepsilon}
\newcommand{\dd}{\,\mathrm{d}}
\begin{document}
\section*{UnsolvedMath ID: 30000081}
\textbf{Source problem code:} catalogue entry
\subsection{Definitions and mathematical statement}
\textbf{Attributed standard-matching conjecture.} Let \(k\) be a field, \(S=k[x_1,\ldots,x_n]\), \(\mathfrak a\subset S\) a monomial ideal with minimal monomial generating set \(G(\mathfrak a)\), and \(A=S/\mathfrak a\). Suppose \(M\) is a \emph{standard matching} on the Taylor resolution. The source asks whether the multigraded Poincar\'e series is\[P_A(\mathbf x,t)\stackrel{?}{=}\frac{\prod_{i=1}^n(1+tx_i)}{1+\sum_{I\in\mathcal U(M)}(-1)^{\operatorname{cl}(I)}m_I\,t^{\operatorname{cl}(I)+|I|}},\] where the summation is over subsets of generators left unmatched by \(M\), \(m_I=\operatorname{lcm}\{u:u\in I\}\), and \(\operatorname{cl}(I)\) is the number of connected components of the graph with vertex set \(I\) and an edge between distinct \(u,v\) when \(\gcd(u,v)\ne1\). Monomials \(m_I\) are read as monomials in the formal multigrading variables \(\mathbf x\), and \(|I|\) is the number of generators in \(I\).\par Under the usual augmented graded-algebra convention, \(P_A(\mathbf x,t)=\sum_{j\ge0}\sum_{\alpha\in\N^n}\dim_k\operatorname{Tor}^A_j(k,k)_\alpha\,\mathbf x^\alpha t^j\), where \(k\) is the residue field at the homogeneous maximal ideal. This convention requires a proper positive-degree monomial ideal, and use of the numerator with all \(n\) variables also needs the intended minimal-presentation hypotheses checked, especially if \(\mathfrak a\) contains linear generators. A frequently relevant restriction is \(\mathfrak a\subseteq(x_1,\ldots,x_n)^2\), but it is not stated in the record and is not silently added as a source hypothesis.\par The definition of \emph{standard matching}, including its additional conditions and which cells survive, is not supplied. Taylor cells are indexed by subsets of \(G(\mathfrak a)\), but this does not define the special matching class. The empty-cell convention also matters: if \(\varnothing\) were included in the displayed sum with \(m_\varnothing=1\) and \(\operatorname{cl}(\varnothing)=0\), the denominator would have constant term \(2\), inconsistent with the usual constant term \(1\) of \(P_A\). An exclusion of the empty subset or another explicitly justified reduced convention is therefore needed; the notation \(\mathcal U(M)\) in the displayed source formula is retained subject to this unresolved qualification. It cannot simply mean every unmatched subset under arbitrary conventions.\par The original prose attributes the conjecture to Charalambous--Reeve and refers to the standard-matching product formula; the title instead spells Charalambous--Reeves. The dated literature triage checked 2026-08-21 labels the record open but explicitly says its extraction is incomplete and openness cannot be verified. The later curated field supplies the formula and carries a \texttt{corrected\_verified} label reporting source inspection; these layers do not supply the missing matching definition or reconcile the dated warning.
\subsection{Short English statement}
The conjecture proposes a rational formula for the multigraded residue-field Poincar\'e series using unmatched Taylor cells. The special matching conditions, empty-cell convention, and minimal-presentation assumptions are missing, so the displayed formula is retained with those qualifications rather than asserted for arbitrary matchings.
\subsection{Sources}
\begin{itemize}\item Local UnsolvedMath registry, permanent ID 30000081, code catalogue entry; title \emph{Charalambous--Reeves Standard-Matching Conjecture}, while original prose says Charalambous--Reeve; Oberwolfach Reports Open Problems.\item \emph{Combinatorial Commutative Algebra (2005)}, Oberwolfach Reports, DOI 10.4171/owr/2004/32, pp. 1703--1746. \url{https://doi.org/10.4171/owr/2004/32}. Proposed year and year tag: 2004; no separate proposer field supplied.\item Supplied literature label \emph{Original Oberwolfach report}: \url{https://doi.org/10.4171/owr/2004/32}. No separate authoritative conjecture citation or definition of standard matching is supplied.\item Literature triage checked 2026-08-21, status \texttt{open}, expressly qualified as not verifiable from its fragment; record updated 2026-09-16. The \texttt{corrected\_verified} field reports checking the official report, but that check is not independently verified here.\end{itemize}
\end{document}
